\scanpage{33}
\uline{NB} Deux filtrations \deleted{différentes}\inserted{équisingulières} de $X$ peuvent très bien définir des \deleted{filtrations} partitions identiques de $X$ en feuilles. C'est la partition en feuilles qui est le plus important, non la filtration.

P.ex. dans le contexte ``par morceaux'', où on dispose donc d'une notion de ``dimension'' et de ``codimension'', il y a au moins 4 façons canoniques, \deleted{\unsure{deux à deux}}\inserted{distinctes} ``en général'', pour décrire une partition équisingulière donnée à l'aide de filtrations : 1°) $X_i =$ réunion des feuilles de dimension $\leq i$, 2°) $X_0 =$ réunion des feuilles minimales $\ldots$ $X_i =$ réunion des feuilles minimales parmi celles qui sont $\notin X_{i-1}$, 3° et 4° définitions ``duales'' de 1° et 2°, donnant des filtrations \uline{décroissantes}.

Pour la suite, on préférera utiliser le langage des espaces à partitions équisingulières, plutôt que des espaces filtrés équisinguliers.

\inserted{\uline{Prop.} Si la partition donnée de $X$ est équisingulière, \unsure{de $=$} pour ttes les parties loc. fermées de $X$ qui sont réunions de feuilles, \deleted{\unsure{la partition}}\inserted{\unsure{induite}}\deleted{\unsure{est}} équisingulière.}

Soit $(X_i)_{i\in I}$ un diagramme d'espaces, et supposons que tout $X_i$ est muni d'une partition \deleted{admissible}. On \deleted{\unsure{en définit le}}\inserted{\unsure{dit qu'un|dit que ce}} système \deleted{\unsure{est dit}} de partitions est \uline{admissible} si

a) $\forall i\in I$, la partition donnée de $X_i$ est admissible

b) $\forall$ flèche $i\to j$, l'application $X_i\unsure{-|\to}X_j$ applique une feuille de $X_i$ dans une feuille (évidemment uniquement déterminée) de $X_j$.
\mspage{34}
